\hwhat{Each non-holdout village day (279) we computed three fluctuation statistics among agents: multi-information, susceptibility (variance of the collective mode) and a heat-capacity analogue (variance of the alignment energy). They were computed for 1-min activity spins, 4-state behavior and 30-min content vectors, as excesses over shuffled surrogates. The pre-registered alarm fires when their trailing z-score reaches 2. \textbf{Real data, mixed:} the alarm caught 10 of 33 goal changes within a day, at a 10\% placebo-window false-alarm rate (AUC 0.68, random-date $p=0.03$). The signal is entirely content: activity is at chance once roster and day-length changes are excluded, and $\sim$40\% of the content signal is a within-day drift. The plain day-to-day shift of the mean content vector scores AUC 0.95. Scaffold changes: 1/13. \textbf{Round 1b} (fixed activity table, restatements removed, two embedding models): hits fall to 0.21--0.27 and the random-date $p$ rises to 0.08 (bge) / 0.21 (gte), at the failure boundary; content (AUC 0.74--0.78) and topic shift (0.95--0.97) replicate; NE40 and NE45 go unseen. Holdout not run.}
\hmodel{Models 01 (inverse Ising) and 11 (vector spins). $s_i(t)=\pm1$: agent $i$ active in minute $t$; $M=\sum_i s_i$; $\hat{\mathbf s}_i(w)$: whitened, unit content vector of agent $i$ in window $w$. $\mathrm{VR}$: susceptibility as a variance ratio (1 for independent agents); $c$: heat-capacity analogue with the uniform alignment energy $E$; $I_G$: Gaussian multi-information ($R$ = correlation matrix); $e(w)$: content alignment energy. Every statistic $X$ enters as observed minus surrogate mean, then a trailing z over the previous 10 days.}
\hmath{\vspace{-8pt}
\begin{gather*}
\mathrm{VR}=\frac{\mathrm{Var}_t M}{\sum_i \mathrm{Var}_t s_i},\qquad c=\frac{\mathrm{Var}_t E}{n},\quad E=-\frac{1}{n}\sum_{i<j}s_is_j\\
I_G=-\tfrac12\log\det R,\qquad e(w)=-\big\langle \hat{\mathbf s}_i(w)\cdot\hat{\mathbf s}_j(w)\big\rangle_{i<j}\\
z_X(d)=\frac{X(d)-\mathrm{med}_{10}X}{\mathrm{sd}^{\rm trim}_{10}X/c_{10}},\qquad Z_{\rm phys}=\big\langle Z_I,Z_\chi,Z_C\big\rangle\ \ge 2
\end{gather*}\vspace{-6pt}}
\hobs{../figures/summary_obs.pdf}{(a) Mean alarm scores around 33 goal kickoffs. The physics alarm (blue) and its content members (orange) rise on the kickoff day but stay below threshold on average. The activity members (aqua) barely move. The content-centroid shift (yellow) jumps. (b) Hit rate against placebo-window false-alarm rate as the threshold is swept; markers sit at threshold 2. The star is the post hoc rule (R1~$\ge$~3 or $Z_{\rm cont}\ge2$): 58\% hits at 2\% false alarms, in-sample.}
\htelos{As posed, a weak alarm. Swarm-wide fluctuation statistics do not light up when an LLM swarm reorganizes. Activity fluctuations react to who is present and for how long; content fluctuations react mostly to the within-day settling after a new goal. A first-moment detector, ``did the topic move since yesterday?'', is cheaper and much better. The physics adds one usable thing, pending the holdout: the content-fluctuation score is a good false-alarm filter beside the topic-shift detector. It matches the shift detector's hit rate at about a tenth of the false alarms. Scaffold changes, the events an operator most wants flagged, stay invisible to all of these daily statistics. An honest negative for ``susceptibility peaks at transitions''.}
\hbottom{Multi-information and susceptibility of an LLM swarm barely mark its reorganizations: on corrected data a quarter of goal changes are caught, all through content, no better than random dates in one of two models, and a plain topic-shift detector does far better.}
\hresults{\footnotesize\noindent\setlength{\tabcolsep}{2pt}\begin{tabular}{@{}p{0.30\columnwidth}p{0.52\columnwidth}p{0.14\columnwidth}@{}}
\textbf{Prediction} & \textbf{Round 1 $\to$ 1b bge / gte} & \textbf{Verdict}\\\hline
P1 goal hit $\le$40\%, AUC $<$0.65 & hit 0.30 $\to$ 0.27 / 0.21; AUC 0.68 $\to$ 0.68 / 0.66; random-date $p$ 0.03 $\to$ 0.08 / 0.21 & mixed $\to$ mixed / failed\\
P2 content $>$ activity & AUC 0.77 vs 0.55 $\to$ 0.78 / 0.74 vs 0.54 (trimmed 0.50) & held\\
P3 rival R1 wins & R1 0.95 $\to$ 0.95 / 0.97; $\Delta$AUC $-0.26\to-0.29$ / $-0.34$ & held\\
P5 scaffold at chance & 1/13 $\to$ 1/13 / 2/13 & held\\
P6 stalls cause false alarms & 0 $\to$ 2/2 false alarms on top stall days (mask doesn't help) & failed $\to$ weak\\
P7 per-day FAR 0.03--0.12 & 0.033 $\to$ 0.036 / 0.036 & held\\
C3 rule R1$\ge$3 or $Z_{\rm cont}\ge2$ & 0.58 hits, 0.02 FAR $\to$ 0.55, 0.02 / 0.67, 0.13 & model-dep.\\
NE39 chat closed (native) & content silent; activity alarm day +1 ($N=4$) & failed\\
NE40, NE45 undocumented (native) & silent, as predicted: not detected blind & held\\
NE43b nudger off (native) & $Z_{\rm act}$ 3.1, R1 2.2 on day 0 (predicted silent) & failed\\
\end{tabular}}
\hobsb{../figures/synthetic_col.pdf}{Synthetic validation at village sampling (30 runs per scenario; AUC of the switch window against placebo windows). Coupling switches in content are recovered; in activity only partly; day-boundary field steps and goal switches are invisible to the fluctuation statistics but obvious to the first-moment rival R1.}
\hcaveats{Daily resolution only. 55 placebo days in 1b (61 in round 1; the new NE dates exclude six), none in regime II. Kickoffs are mostly Mondays. The verdict sits on the random-date boundary and moves with the embedding model and the dedupe (bge without dedupe: $p=0.11$). The activity channel reacts to roster, day length and, on the fixed table, joint silences. The operator rule was chosen post hoc and its false-alarm rate is model-dependent. The NE43b alarm is a single, unreplicated day. Synthetic figure (left) is round 1's generator.}
\hnext{Re-freeze \texttt{confirm.py} on the fixed table (both models) before any holdout run. Build HH268's combined monitor: topic shift (R1) plus H56's log features and H38's schedule signals; intraday windows.}
